%--------------------------------------------------------------------------------
%
%
%--------------------------------------------------------------------------------
\pagebreak
\section*{LD — Load Register}
\addcontentsline{toc}{section}{LD — Load Register}

\begin{iPageDescEntry}{Syntax}
    \texttt{LD[.D/W/H/B] RegR, SignedImmed13( RegB )} \\
    \texttt{LD[.D/W/H/B] RegR, RegX( RegB )}
\end{iPageDescEntry}

\begin{iPageDescEntry}{Format}
    The \texttt{LD} instruction uses the indexed and register-indexed instruction formats. \\

    \begin{flushleft}
    \drawInstrFormat
        {0,1,3,5,5.5,6,6.5,8.5,9.5,16}
        {\OpCodeGrpMEM, \OpCodeLD, RegR, 0, M, 0, RegB, dw, S-IMM-13}
    \end{flushleft}

    \begin{flushleft}
    \drawInstrFormat
        {0,1,3,5,5.5,6,6.5,8.5,9.5,11.5,16}
        {\OpCodeGrpMEM, \OpCodeLD, RegR, 0, M, 1, RegB, dw, RegX, 0}
    \end{flushleft}
\end{iPageDescEntry}

\begin{iPageDescEntry}{Description}
    The \texttt{LD} instruction loads a value from memory into a general-purpose register.
    - The indexed format adds a signed immediate to \texttt{RegB} to form the effective address.
    - The register-indexed format adds \texttt{RegX} to \texttt{RegB} to form the effective address.

    The instruction may trigger a memory reference trap if the address is invalid.
\end{iPageDescEntry}

\begin{iPageDescEntry}{Operation}
    \texttt{RegR <- memLoad( memOperand() );} \quad (for indexed formats)
\end{iPageDescEntry}

\begin{iPageDescEntry}{Exceptions}
    - Data Alignment Trap \\
    - Memory Access Violation Trap \\
    - Memory Protection Violation Trap \\
    - Data TLB Miss Trap
\end{iPageDescEntry}

\begin{iPageDescEntry}{Notes}
    - The instruction supports different sizes: byte (.B), halfword (.H), word (.W), doubleword (.D). \\
    - Often used in loops to fetch data from memory.
\end{iPageDescEntry}


